\documentclass[11pt]{article}
\usepackage[margin=1in]{geometry}
\usepackage{amsmath,amssymb,booktabs,graphicx}
\usepackage[hidelinks]{hyperref}
\usepackage[T1]{fontenc}
\makeatletter
\let\tableinput\@@input
\makeatother
\newcommand{\Prs}{\widehat P}
\title{Can a busier market be said to have more persistent order flow?\\[3pt]\large Sharp bounds from anonymous execution records}
\author{Yunjie Guo \and Hanqing Yang\thanks{Corresponding author. Email: hanqing.yang189@gmail.com. Both authors are independent researchers in Suzhou, China.}}
\date{}
\begin{document}
\maketitle
\begin{abstract}
A marketable order can consume several resting orders at several prices. This leaves the denominator of order-level trading statistics uncertain even when the aggressor side of every execution is known. We derive the exact finite set of one-step same-side frequencies consistent with an anonymous execution tape under side, price-priority and timing restrictions. A linear scan gives attainable endpoints even when a sweep can span timestamps up to a specified duration. For comparisons across activity states we also calculate the fewest record joins needed to erase a reported ranking. For 8.03 million Binance Spot records across 12 asset-days, both fill and aggregate-record counts rank the busiest hours higher on every day. Five rankings survive every native-timestamp monotone reconstruction. Two of the seven reversible rankings require the busier state's fraction of possible joins to exceed the quieter state's by over 98 percentage points. A separate Coinbase audit with public taker identifiers recovers the labelled episode count by maximal native-time joining in 23 of 24 windows. The exception identifies a limit of transferring sweep assumptions across venues.
\end{abstract}

\section{The measurement problem}
Suppose a market monitor reports that buy orders follow buy orders and sell orders follow sell orders more frequently in its busiest hours. Counts of elementary fills give greater weight to orders that encounter many resting counterparties. A feed aggregating fills at the same execution price removes part of this weighting. An incoming order walking across several prices can still generate several aggregate records. If order size or displayed depth changes with activity, the effect on a comparison can differ across the two states.

This paper studies the fraction of adjacent aggressor orders with the same sign, pooling buy and sell pairs. The target concerns immediately executed portions of incoming orders. A trader can send many such orders while executing a larger investment decision. We make no inference about investor identity, metaorder length, or the decay of long-lag correlations. In particular the baseline under independent signs with buy probability $q$ is $q^2+(1-q)^2$. Differences in buy share can change the statistic even in the absence of serial dependence.

Prior work establishes both the importance of choosing the event clock and the hazards of assigning fills to orders. T\'oth et al.\ [1] separate splitting from herding using member identifiers. Schwenk-Nebbe and Thimsen [2] examine how exchange timestamps can link executions in US equity data. Dyhrberg, Foley and Svec [3, online Appendix A2] use direction, monotone execution prices and a time window to join prints. Foley et al.\ [4] study exchange-specific timestamp precision and slow matching. We take these reconstruction restrictions as inputs. Our contribution is to propagate every permitted segmentation into a sharp state contrast, with a certificate counting the joins necessary for a different ranking. This answers a question that benchmarking one reconstruction against another leaves open.

\section{The feasible set}
Within each uninterrupted observation segment the ordered records are $(s_i,p_i,t_i)$ with $s_i\in\{-1,1\}$, exact decimal price $p_i$, and nondecreasing exchange time $t_i$. An atom is one execution record known to belong to a single taker. The latent order tape partitions the atoms into contiguous blocks. All records in a block must have the same sign. A buy sweep moves up the book and a sell sweep moves down. In the primary model every adjacent pair in a block must satisfy $s_i(p_{i+1}-p_i)>0$ and have equal native timestamps. The strict inequality reflects aggregation by price for one taker and stamp. We assume a sweep completes at one stamp in this specification. The data dictionary by itself does not guarantee that timing assumption.

We also consider a strict-price rule at a coarser timestamp, and a maximum elapsed time from the first to the last atom in an order. Under the latter we permit an unchanged price at distinct native stamps. These restrictions are local to each block and hereditary: any contiguous sub-block of a feasible block is feasible. We allow any partition satisfying them. Sharpness refers to this information model. Additional book or order messages could exclude some feasible partitions.

\paragraph{Proposition 1.}
Fix $J$ nonempty segments, $M$ records and $S$ observed sign switches within segments. In each segment repeatedly take the longest feasible prefix as one order. Let $G$ denote the resulting total number of orders. Every order count $n\in\{G,\ldots,M\}$ is attainable. For $n>J$ the exact feasible set of same-side frequencies consists of
\begin{equation}
       \Prs(n)=1-\frac{S}{n-J}. \label{eq:p}
\end{equation}
In particular, when $G>J$ the lower and upper endpoints use $G$ and $M$ respectively.

\paragraph{Proof.}
After $k$ blocks the greedy prefix reaches at least as far as any other feasible $k$-block prefix. For the induction, if the previous greedy endpoint lies inside the competitor's next block, the remaining suffix of that block is feasible. If it already reaches beyond that block the assertion holds immediately. Thus the greedy partition minimizes count in each segment. Splitting any nonsingleton block retains feasibility and increases count by one, up to the singleton partition. A sign change always forces a cut. Every within-block join deletes exactly one same-side transition, leaving $S$ unchanged. There are $n-J$ adjacent order pairs. This proves (\ref{eq:p}). \hfill$\square$

The buy and sell counts also independently attain every integer between their side-specific greedy and singleton counts. This follows by applying the argument within each maximal same-side run. Thus one can bound the order-count buy share or condition on a known buy count without enumerating partitions. In the duration model a chain of individually close arrivals can exceed the permissible total span. Measuring from the current block's first stamp is essential. For fixed rules a single scan uses constant working memory per segment.

\paragraph{A five-record example.}
Take four buys at prices 100, 101, 102 and 103 at one stamp, followed by a sell. The buy side admits between one and four orders. One sign switch gives feasible frequencies $0,1/2,2/3,3/4$. Even an exact timestamp combined with a strict book walk leaves the individual order boundaries unresolved.

\paragraph{State comparisons and a tipping count.}
Let $H$ and $L$ denote disjoint groups of complete segments. Their endpoint counts are independently attainable. The difference $\Prs_H-\Prs_L$ therefore has sharp endpoints $\Prs_H(G_H)-\Prs_L(M_L)$ and $\Prs_H(M_H)-\Prs_L(G_L)$. The interval displays the extrema of a finite set. For a strictly larger record-level statistic in group $A$ than $B$, write $D_j=M_j-J_j$ and suppose both switch counts are positive. The minimum number of joins in $A$ that can erase the ranking, holding $B$ at singletons, is
\begin{equation}
 k^*=\left\lceil D_A-\frac{S_A D_B}{S_B}\right\rceil.\label{eq:k}
\end{equation}
This follows by solving $S_A/(D_A-k)\geq S_B/D_B$. Joining records in $B$ makes the inequality harder. If $k^*>M_A-G_A$ the ranking survives every permitted segmentation. Otherwise $k^*$ is attainable (with nonzero denominators). The ratio to available joins gives a sensitivity measure without assigning probabilities to partitions.

A second sensitivity calculation allows joins in both states. Write $U_j=M_j-G_j$ for merge capacity and $k_j$ for actual joins. For $U_A,U_B>0$ measure the extra concentration of joins in the initially winning group by $[k_A/U_A-k_B/U_B]_+$, where $[x]_+=\max(x,0)$. Given $k_B$, the smallest count that erases the ranking is
\[
 k_A(k_B)=\max\!\left(0,\left\lceil D_A- S_A(D_B-k_B)/S_B\right\rceil\right).
\]
Scan integers $0\leq k_B\leq U_B$, discarding $k_A(k_B)>U_A$. Minimizing the positive excess over the remaining pairs gives the exact smallest excess needed. Larger $k_A$ at fixed $k_B$ cannot improve the objective. Proposition 1 supplies a partition for each witness pair. A zero excess means that the winner need not join a larger proportion than the loser. It does not require exactly equal integer fractions. This secondary diagnostic was developed after the initial labelled audit.

\section{Markets, observations and activity states}
We use Binance Spot's daily aggregate-trade archives [5,6] for ETH/USDT, SOL/USDT and LTC/USDT. For each pair we take the fifteenth day of January, April, July and October 2025. The design fixes one full UTC day per quarter during the first year for which the archive documents microsecond spot timestamps. Each row gives a price, buyer-maker flag and first and last elementary trade IDs. A buyer acting as maker implies an aggressive sell. The inclusive ID span supplies the number of elementary fills. Within each file we check the publisher's SHA-256, record ordering, ID ordering and timestamp ordering. We break segments at each UTC hour and at any gap between successive ID ranges or aggregate IDs. All event-clock comparisons use these same breaks.

Within each symbol-day, rank the 24 hours by elementary fill count. Use the six highest and six lowest hours as the high- and low-activity groups, breaking ties by hour. This assignment is observed before reconstructing any orders. Each group pools transition counts across its hours. Hours with more transitions consequently receive more weight. The labels describe contemporaneous trading intensity and carry no ex ante or causal interpretation. The other twelve hours contribute to coverage figures only. We compare exact extrema with two plug-ins, one event per elementary fill and one per aggregate record. The former can lie beyond the order-level bounds because the feed already links fills within an aggregate. The bounds are deterministic conditional on the sampled tapes. Trade-level standard errors would address a different source of uncertainty.

\section{Initial Binance activity comparisons}
The twelve archives contain 27,640,107 elementary fills in 8,026,298 aggregate records. All 288 hours have observations, and the ID checks find no gaps. The three point reconstructions agree on the direction in every symbol-day. High activity yields more successive same-side events when one counts fills, counts aggregate records, or maximally joins records in both states. Yet the sharp native-time contrasts in Table~\ref{tab:primary} straddle zero in seven comparisons. ETH supplies four of these. All four ETH record contrasts are positive and below 2.6 percentage points, while even the narrowest ETH interval exceeds 18 points in width.

\begin{table}[htbp]\centering
\caption{Difference in same-side frequency, busiest minus quietest hours. All entries after month are in percentage points.}\label{tab:primary}
\small\begin{tabular}{llrrrrrr}\toprule
Asset & Month & Fills & Records & Max joins & Lower & Upper & Excess\\\midrule
\tableinput primary.tex
\bottomrule\end{tabular}

\smallskip\begin{minipage}{0.94\linewidth}\footnotesize
Max joins applies the native-time greedy partition in both groups. Lower and upper optimize the partitions separately. Excess is the minimum positive difference in the fractions of available joins (initial winner minus loser) required to erase the record ranking. A dash denotes an unreachable tie or reversal. Endpoints use unrounded counts.\end{minipage}
\end{table}

The two April and October SOL comparisons illustrate the difference between a formal reversal and a plausible one. In April the high-activity group needs at least 5,197 of its 5,290 available joins if the quieter group uses none. In October it needs at least 25,272 of 25,312. Even allowing both states to join, minimum excesses are 98.17 and 99.80 percentage points. A reconstruction joining comparable fractions in both states therefore preserves both rankings. By contrast, for July ETH a difference of 12.56 points suffices. The repository supplies the integer witness counts giving these minima.

\paragraph{How many exceptions could eliminate a robust ranking?}
For the native rule, joinability is a property of each adjacent edge. Keep side changes and segment breaks mandatory, but allow up to $b$ joins that fail the native time or price test. If one state has $U=M-G$ eligible joins, its capacity then rises to $U+\min(b,G-J-S)$: the $G-J-S$ remaining within-segment edges connect equal signs and can be joined independently in this relaxation. To erase a positive ranking at minimum cost, put all exceptions in the winning state, again leaving the loser at singletons. The exact additional budget needed is $[k^*-U_A]_+$. For April LTC, 227 exceptional joins suffice beyond its 2,076 eligible ones. For January and July SOL, the corresponding budgets are 27,158 and 29,600. This measures tolerance to misspecification in boundary counts without extrapolating an error rate from another exchange.

\paragraph{Sensitivity to time and direction.}
Coarsening the timestamps to millisecond buckets while retaining strict price progression increases median contrast width from 4.19 to 4.66 percentage points (Table~\ref{tab:timing}). The 20 ms convention additionally admits some flat-price edges and lets a sweep cross bucket boundaries. It widens every interval enough to permit either ordering. Its contrast with the primary result reflects both timing and price assumptions.

\begin{table}[htbp]\centering
\caption{Width of the sharp state contrast across 12 symbol-days (percentage points) and number admitting both signs.}\label{tab:timing}
\begin{tabular}{lrrrr}\toprule
Restriction & Min & Median & Max & Both signs\\\midrule
\tableinput timing.tex
\bottomrule\end{tabular}
\end{table}

As a secondary check, weight buy-conditioned and sell-conditioned continuation equally. For predecessor sign $s$, let $F_s$ count switches away from $s$ and $E_s$ count segments ending in $s$. Writing $n_s$ for the feasible order count of that sign gives the balanced statistic
\[
  1-\tfrac12\sum_{s\in\{+,-\}} F_s/(n_s-E_s).
\]
The numerator and terminal signs are invariant to joining. Both denominators are positive in the sampled groups. By side-specific attainability, the lower and upper endpoints follow by taking the two side minima and the two maxima. Five of twelve native-time comparisons still have positive lower endpoints, but the membership changes. The April SOL lower endpoint rises from $-0.026$ to $+0.031$ percentage points. The April LTC lower endpoint falls from $+0.241$ to $-2.074$ points. Thus the predecessor-side mix affects which activity rankings are identified even when the total number of surviving comparisons stays fixed.

\section{Short-window check of the sweep assumptions}
Coinbase Exchange match messages identify the taker order [7]. We replay archived match messages [8] for BTC-USD, ETH-USD and SOL-USD over the first ten minutes of UTC trading on the first day of January, April, July and October 2025. This fixed archive sample replaced an unavailable live WebSocket connection. After inspecting these 12 product-windows we repeated the audit on the first day of the following month in each quarter, at 12:00 UTC. The two panels contain 56,657 fills overall.

To reproduce the observation structure, we first combine adjacent fills with identical taker ID, execution price, side and timestamp. Labels are needed for this within-level aggregation, just as the Binance feed producer uses order information. We then hide the labels and reconstruct across the resulting records. Thus this exercise tests the unknown \emph{between-record} boundaries. Ground truth for this contiguous-partition model is a consecutive episode of the same taker ID. The 24 streams have no trade-ID gaps. In three windows an ID reappears in a later episode, illustrating the distinction between contiguous execution episodes and distinct order IDs. The totals contain 33,060 episodes and 33,057 within-window distinct IDs.

\begin{table}[htbp]\centering\small
\caption{Labelled episodes and greedy reconstructed counts on Coinbase. Each row pools four ten-minute windows.}\label{tab:validation}
\begin{tabular}{llrrrrrr}\toprule
UTC & Product & Fills & Atoms & Label & Native & 1 ms & 20 ms\\\midrule
\tableinput validation.tex
\bottomrule\end{tabular}
\end{table}

The labelled partition obeys the native rule and achieves its lower endpoint in each of the twelve midnight windows and eleven of the twelve noon windows. For this edge-local rule maximal joining has a unique minimum partition, so agreement in count plus feasibility implies agreement in boundaries. The single exception occurs in ETH-USD in August. Two consecutive sell-side records at USD 3,656.84 and 3,656.60 have the same taker identifier but are 1,077,664 microseconds apart. The direction and price progression are admissible. The elapsed time exceeds even the 20 ms allowance.
For this window even the direction of the endpoint error matters: the native rule requires at least 2,213 orders but the labelled tape has 2,212 episodes. Under the coarser rules the count falls inside the feasible range although the labelled partition still violates one edge constraint. Counting coverage alone would miss this failure.

The labelled venue provides two lessons for the anonymous venue. First, the regime that joins nearly all eligible native-time boundaries in \emph{both} states deserves attention as an empirical benchmark. The narrow SOL tipping margins quantify how far one would have to depart from that behavior to erase a ranking. Second, admissibility of the true partition is an empirical assumption. Treating Coinbase labels as Binance labels would require evidence about the two engines and feed construction that these data cannot supply.


\section{Scope and implications}
Identification of a state ranking requires knowing enough about within-order fragmentation on each side of the comparison. Adding more observations at the same information level need not determine that fragmentation. Each ambiguous boundary contributes another feasible choice. A study with parent identifiers can instead calculate where the observed partition sits in the feasible set and test which restrictions it violates.

The analysis uses a small deterministic calendar sample, a common quote currency and one primary venue. It does not estimate a population frequency of unidentified rankings. The two activity groups can also have different buy shares, order sizes, liquidity and news exposure. The extrema permit different joining behavior in the two states. A model imposing a common probabilistic rule would narrow the admissible set at the cost of additional assumptions. Extending the statistic to several order lags requires tracking the positions of switches, since their total alone ceases to determine agreement. For the one-step question a report of the two sharp endpoints and the tipping count makes explicit exactly which ordering the anonymous tape can support.

\paragraph{Replication and disclosures.}
The project repository contains the protocol, data retrieval and integrity checks, exhaustive small-tape tests, sufficient statistics and table-building code:
\url{https://github.com/fightforbelief/p14-order-flow-bounds}. The authors report no specific funding and no conflicts of interest.

\begin{thebibliography}{9}
\bibitem{toth} B. T\'oth, I. Palit, F. Lillo and J. D. Farmer. Why is equity order flow so persistent? \emph{Journal of Economic Dynamics and Control} 51 (2015), 218--239.
\bibitem{taq} S. Schwenk-Nebbe and C. Thimsen. The participant timestamp: Get the most out of TAQ data. Working paper, \href{https://ssrn.com/abstract=3744743}{SSRN 3744743}.
\bibitem{tick} A. H. Dyhrberg, S. Foley and J. Svec. When bigger is better: The impact of a tiny tick size on undercutting behavior. \emph{Journal of Financial and Quantitative Analysis} 58 (2023), 2387--2416.
\bibitem{slow} S. Foley, W. Krekel, V. Mollica and J. Svec. Not so fast: Identifying and remediating slow and imprecise cryptocurrency exchange data. \emph{Finance Research Letters} 51 (2023), 103401.
\bibitem{data} Binance. \href{https://github.com/binance/binance-public-data}{Public market data archive and schema}.
\bibitem{api} Binance. \href{https://developers.binance.com/docs/binance-spot-api-docs/rest-api/market-data-endpoints}{Spot market data endpoints, compressed/aggregate trades list}.
\bibitem{cb} Coinbase. \href{https://docs.cdp.coinbase.com/exchange/websocket-feed/channels}{Exchange WebSocket channels, match message}.
\bibitem{tardis} Tardis.dev. \href{https://docs.tardis.dev/api/http-api-reference}{Historical raw data feeds, HTTP API reference}.
\end{thebibliography}
\end{document}
